\section{Discussion}
\label{sec:discussion}

\subsection{Principal Findings, Lesion Fidelity Trade-Offs and Capacity Controls}
\label{sec:d-principal}
Several pre-specified results did not hold. The soft-tissue image alone did not improve any downstream task
(Tables~\ref{tab:downstream_cls} and~\ref{tab:downstream_det}); a multi-stream classifier that read the radiograph, soft-tissue and bone
images did not improve on the radiograph alone (H6, Table~\ref{tab:cls_primary}); SSIM fell short of the reference model by more than its
margin (Table~\ref{tab:h2}); and conformal intervals calibrated on synthetic data covered 14\,\% instead of 90\,\% of real bone pixels
(Table~\ref{tab:reliability}), 89\,\% after an exploratory recalibration on 24 real pairs (Table~\ref{tab:bone_recal}). Real-nodule contrast
was higher than with the reference model on JSRT and lower on Node21 (Table~\ref{tab:nodule_cnr}). In Study~III, distillation from CT did
not improve real-radiograph classification (H11, H12), the CT and X-ray Grad-CAM++ maps did not agree (H14), and the Swin classifier trailed
DenseNet-121 over the 14 diseases (H13), almost entirely through one rare class and with a reference trained twice as long at four times the
encoder learning rate (Table~\ref{tab:xm_hyp}).

The central results of Study~I held. Lesion absorption fell steadily with the twin weight (Table~\ref{tab:dose_response}); on an audit set
no model had seen during development, the twin ensemble absorbed a third of what the released CT-supervised model absorbed and less than a
third of what the matched control absorbed (Table~\ref{tab:audit_contrasts}), at non-inferior bone-region fidelity; and the effect
replicated with two other encoders (Table~\ref{tab:encoders}) and on paediatric radiographs (Section~\ref{sec:res-bdcxr}).

\emph{What the audit adds.} Fidelity metrics and lesion absorption rank suppressors differently (Fig.~\ref{fig:pareto}). The v1 model has
the highest bone-region PSNR of the CT-supervised models and the highest absorption; ours and the reference model are within 0.06\,dB of
each other and a factor of three apart in absorption. BS-LDM, trained on dual-energy pairs from the evaluation hospital, is about 7\,dB ahead
in fidelity on that hospital's images, yet absorbed 0.42 of an inserted lesion in the development audit (at unpaired sites, so only
descriptive): fidelity to a dual-energy target does not protect what the target never showed. Part of that lead is the appearance of the
site's dual-energy images rather than suppression: a smooth calibration learned from 24 pairs of the site raised every CT-supervised model
by 2.1 to 2.4\,dB in the bone region without raising lesion absorption, and BS-LDM remained about 5\,dB ahead
(Section~\ref{sec:res-sitecal}; exploratory). Absorption must still be read beside
suppression strength. The reference model removes more rib-edge energy than ours (48\,\% against 30\,\%) and more non-rib energy (22\,\%
against 8\,\%), so part of its higher absorption may reflect stronger suppression in general.

\emph{What the regulariser costs.} On real paired radiographs the twin term costs no measurable bone-region fidelity up to $w=16$ (21.22 to
21.28\,dB across $w=0\ldots16$); at $w=32$ and 64 PSNR drops by about 0.3\,dB, so $w=16$, chosen on development data before the freeze, is
the usable operating point. The cost appears where the truth is exact: on the Qure-projection test split the bone error of the twin ensemble
is 41\,\% higher than that of the control (MAE 0.0041 against 0.0029). The regulariser thus trades synthetic bone accuracy for lesion
invariance, a trade not visible on real images. The ablation shows that the effect needs gradients through both twin branches and a
consistency term confined to the pasted lesion; showing lesions during training without the twin term (CLIA) did not lower absorption. The
reference model is a SwinV2-B at $1024$~px trained with a far larger budget, so only the matched no-twin control, which absorbs slightly more
than the reference model (0.178 against 0.149), isolates the regulariser. After the exploratory site calibration the no-twin control was
slightly ahead of ours (23.74 against 23.58\,dB), a cost of the regulariser that the uncalibrated comparison did not show; it was not
tested. Diffusion models trained on real dual-energy pairs
\cite{sun2026bsldm,sun2026debonedit} need paired data from the target site, sample slowly (about 2.5 minutes per image for BS-LDM on a T4
GPU) and were not designed for lesion invariance.

\emph{Do classifiers gain from suppressed inputs?} The multi-stream hybrid matched the radiograph alone on all three datasets; on NIH the
interval rules out a gain larger than 0.0024 in macro-AUROC, so the null result is informative. The hybrid also matched the capacity
control, and plain channel stacking did at least as well as separate streams. Eq.~\eqref{eq:info} offers one explanation: the suppressor
adds a transformation learned from CT but no information, and a classifier trained on thousands of radiographs may learn whatever
bone-invariant features its labels reward. Nodules and masses had the largest estimates on both adult datasets ($+0.008$ to $+0.009$ for
the soft-tissue input), and the pointing game agreed: maps of the radiograph stream found focal lesions less often than other findings,
whereas maps of the soft-tissue stream found both about equally often. Two cautions apply. The NIH test split had already been used by the
downstream classifiers of Study~I, so the NIH result is a re-test with new models; and VinBigData, the independent test, has a single merged
nodule/mass class whose interval included zero. The loss for diffuse findings appeared on NIH only, mostly for infiltration; with
report-derived NIH labels \cite{wang2017chestxray} and no such loss on the radiologist-labelled VinBigData, these data cannot tell whether
suppression removes cues for diffuse disease or interacts with label noise.

\subsection{Cross-Modality Transfer Limits, Localization Faithfulness and Robustness}
\label{sec:d-xm}
\emph{Why CT did not teach the radiograph classifier.} The distilled classifier matched the one started from ImageNet within $\pm0.0016$
AUROC on the shared findings and matched the classifier pretrained on the same projections without the teacher. First, the teacher had
little to give: on the six findings it was hardly more accurate on CT (0.821) than $P_0$ on the projections of the same test patients
(0.814), and it was weakest for lung lesions (0.608), where CT should help most. Second, the 2.5D recipe was fragile, with one collapsed and
one late-learning run in seven, so students learned from teachers of uneven quality. Third, gains on the projections had to survive the move
to real radiographs: the zero-shot results (0.640--0.648 on NIH) show how wide that gap is, and ten epochs of fine-tuning on \num{77890}
radiographs may override the initialisation, as the equal results of $F_{\mathrm{p0}}$ and $F_{\mathrm{in}}$ suggest. The negative result
therefore concerns this teacher and this amount of paired CT, not cross-modal distillation in general.

\emph{Explanations and robustness.} The two modalities did not produce maps of the same kind. The teacher's Grad-CAM++ mainly ranks slices,
so its projected map varies from top to bottom but hardly across, while the $8\times8$ token grid of the last SwinV2-T stage gives
Grad-CAM++ maps that spread along the image borders. On real radiographs Swin Grad-CAM++ hit expert boxes less than half as often as
DenseNet-121 Grad-CAM++, whereas Integrated Gradients and decoder attention localised about as well for both encoders; agreement between
modalities should therefore be measured with maps that localise. Robustness, in turn, was not visible in accuracy: the Swin classifiers were
not more accurate than DenseNet-121, but noise flipped about a third as many finding-free images (0.27--0.29 against 0.78). Under rib
emphasis the distilled classifier flipped fewer images than the one started from ImageNet ($p=0.044$, uncorrected), but so did the classifier
pretrained on projections without the teacher, which points to pretraining rather than to CT.

\emph{Uncertainty.} The bone-image intervals held where they were calibrated and failed under the shift to real radiographs; 24 real pairs
from the target site sufficed to restore coverage here, although transfer between sites is untested. For the classifiers, conformal risk
control kept its guarantee for every input and encoder, but label sets holding about half of the 14 NIH labels are of little use to a
reader, and the suppressor's own uncertainty made abstention worse in every configuration: uncertainty about the bone image is not
uncertainty about the diagnosis.

\subsection{Clinical Translation, Limitations and Future Research}
\label{sec:limitations}
\emph{Recommendations.} A bone suppressor should be audited for lesion absorption, reported next to a suppression measure, before it is
used; suppressed images should be shown beside the original rather than instead of it; and the ensemble spread should serve as a warning
flag rather than a calibrated interval unless recalibrated on real pairs (it rose 2.7-fold on paediatric radiographs). Classifiers should not
be expected to gain more than small amounts of overall accuracy from bone-suppressed inputs, which helped only as extra channels, or from a
CT teacher of this kind, and any benefit should be tested per finding against a control of the same capacity. Swin classifiers of
radiographs should be explained with Integrated Gradients or decoder attention rather than Grad-CAM++, with maps scored against expert
boxes, and claims of robustness should rest on perturbations that add no disease. Fidelity comparisons between CT-supervised and
dual-energy-trained suppressors should separate the appearance of the test site from suppression, for example with a calibration on a few
development pairs reported beside the uncorrected numbers.

\emph{Limitations.} Each limitation is stated with its consequence for how the results must be read.
\begin{enumerate}
\item \emph{Synthetic supervision and one source of real pairs.} The bone image of a real radiograph is never observed. Real-image fidelity
 is measured against dual-energy soft-tissue images from one hospital and one acquisition set-up (GE Discovery XR656), which differ from
 the radiograph in tone and fine texture, so raw and band-limited metrics are both reported and calibrated metrics are not absolute.
\item \emph{The audit is linear and additive, and shares its mechanism with the regulariser.} Lesions are added in image units; beam
 hardening, scatter and tone mapping are not modelled, so $A$ measures the leak of an added signal rather than clinical detectability, and
 it is read next to a suppression gain because $A=0$ is reached by doing nothing. The real-nodule contrast and downstream subgroup analyses
 limit, but do not remove, the circularity of inserting lesions both in training and in the audit.
\item \emph{Development use of SZCH.} All 75 SZCH images were used in development (audit) and 24 for paired decisions, so only the 51
 held-out pairs support fidelity claims; the plan changed once after the SZCH pairs had been looked at. The exploratory site calibration
 (SAP\_v5) was designed after H2$'$ was known, re-uses the 51 held-out pairs and is specific to one hospital and scanner, whose own 24
 development pairs it was fitted on.
\item \emph{Resolution, backbone and operating point.} Our suppressors predict at 512\,px with tiny encoders, the reference model at
 1024\,px with SwinV2-B. The SSIM gap sits in fine texture (with a band limit of $\sigma=4$ it shrinks from 0.021 to 0.002); whether
 prediction at 1024\,px closes it was not tested. Only $w=16$ was evaluated with the full protocol.
\item \emph{Uncertainty.} Conformal calibration treats pixels as exchangeable and was calibrated on synthetic data; the recalibration on
 real pairs was chosen after the shift was seen and used 24 pairs from one hospital. Ensemble members share training data.
\item \emph{Scope and external validity.} Training and paired evaluation use adult frontal radiographs; BDCXR-3257 is a single-centre
 paediatric set with image-level labels and no patient identifiers, and its audit inserts adult nodules into children's radiographs. There
 is no lateral imaging and no reader study. On JSRT the bone overlap of each audit site came from the reference model's bone map.
 Downstream models were not tuned per input, and the overlap analysis is not estimable on TBX11K. The models are for research use only.
\item \emph{Classifier study.} One encoder (DenseNet-121) and one resolution ($256\times256$\,px) were tested; heavier encoders or full
 resolution may use the extra streams differently. The NIH test split served the downstream experiments of Study~I as well, its labels are
 report-derived, and VinBigData has one focal class. The per-class DeLong tests treat NIH images as independent although patients
 contribute several images, and the pointing game depends on the attribution method, the 15\,px tolerance and the eight annotated NIH
 classes.
\item \emph{Cross-modality study.} The 2.5D CT recipe was seed-unstable, and the replacement of seed 2 by seed 5 followed a rule fixed
 beforehand but used a teacher trained later. The CT subset is small (\num{1600} training volumes) and its labels are report-derived. The
 localisation of the H13 difference in hernia (86 positive test images) comes from per-class tests planned only as descriptive, and H13
 compares encoders trained with different pre-specified budgets (DenseNet-121 for 20 epochs at the full encoder learning rate, Swin for 10
 epochs at a quarter of it), neither tuned further; the reference's NIH test score was known from Study~II when the fourth plan was
 frozen. H14 relies on Grad-CAM++ for both modalities, the least localising of the three maps on real radiographs. The example heat maps
 contain no mass or nodule, and VinBigData fine-tuning was not run. Teacher outputs, predictions and weights reached the authors' computer
 against the design's own data rule; nothing derived from CT-RATE was released.
\item \emph{Post-freeze corrections and registration.} Analysis errors found in an independent check after the first locked run were
 corrected (among them a resampling error that narrowed interaction intervals and detection metrics that depended on the order of tied
 scores by up to 0.00055); cloud and local GPU outputs differ by at most 0.0024\,dB. No pre-specified verdict changed. The analysis plans
 were hash-frozen in version control before testing but not deposited in a public registry.
\end{enumerate}

\emph{Future research.} Future work should test prediction at $1024$~px, recalibration on real pairs from more sites, other lesion types and
readers, and, for transfer from CT, a stable volumetric teacher trained on many more volumes.
